\section{Discussion}
\label{sec:discussion}

\subsection{Effectiveness of the Evidence Gate}

The 83.2\% relative reduction in TP FNR achieved by the Evidence Gate confirms
that premature investigation termination on empty query returns is the dominant
failure mode in tool-using LLM agents, as identified by Tsegai et al.~\cite{tsegai2026aida}.
The gate introduces a soft pause mechanism: rather than immediately concluding FP,
the agent reissues queries with expanded time windows or alternative host scopes
before triggering human escalation. This behavioral modification requires no
modification to the underlying LLM weights or inference procedure.

\subsection{Operational Cost-Benefit Analysis}

The 18.2\% human escalation rate represents an operational overhead relative to
fully automated baselines. At enterprise scale (e.g., 50,000 daily alerts across
a medium enterprise), this implies approximately 9,100 daily Tier-2 escalations.
Given that standalone GBDT misses 18.4\% of true attacks and standalone LLMs miss
40.4\%, the marginal cost of 18.2\% escalation is offset by the 11.6\% reduction
in missed attacks relative to pure GBDT, translating to an estimated \$2.1 million
in avoided breach costs per 1,000 prevented incidents (derived from IBM 2026
per-incident cost of \$4.99M~\cite{ibm2026databreach} at a conservative 42\%
reduction factor).

\subsection{Threats to Validity}

\textbf{Dataset telemetry anonymization:} Microsoft GUIDE applies privacy-preserving
transformations to customer data, potentially removing low-level process command-line
arguments that constitute important kill-chain indicators. This may systematically
underestimate GBDT feature richness in production deployments where full telemetry
is available.

\textbf{LLM API latency variability:} The reported 1.25-second median latency
assumes stable API endpoints with $<$500ms 95th-percentile call latency. Under
API rate-limiting or provider outage conditions, the system degrades to standalone
GBDT classification with automatic escalation, maintaining safety at the cost of
reduced analytical depth.

\textbf{MITRE ATT\&CK version drift:} The STIX~2.1 knowledge index was built from
MITRE ATT\&CK v16.1. Future adversarial technique evolution may require periodic
index refreshes. The modular Domain Knowledge Engine design enables independent
reindexing without retraining downstream components.
